\subsection{مبانی نظری و فرمول‌بندی ریاضی}
فرض کنید ماتریس وزن پروجکشن اولیه $\mathbf{W}_0 \in \mathbb{R}^{d_{\text{out}} \times d_{\text{in}}}$ با استفاده از تجزیه مقادیر منفرد وزن‌دار فیشر تجزیه شود. در صورتی که $\mathbf{F} \in \mathbb{R}^{d_{\text{out}} \times d_{\text{in}}}$ ماتریس قطری اطلاعات فیشر باشد، تجزیه بر روی ماتریس وزنی $\tilde{\mathbf{W}}_0 = \mathbf{W}_0 \odot \sqrt{\mathbf{F}}$ اعمال می‌گردد:
\begin{equation}
\mathbf{W}_0 \approx \mathbf{U}_{:, :r} \operatorname{diag}(\boldsymbol{\Sigma}_{1:r}) \mathbf{V}_{1:r, :}^\top
\end{equation}
که در آن $r \le \min(d_{\text{in}}, d_{\text{out}})$ رتبه فاکتورسازی است. در لایه \lr{SALTLinear}، وزن موثر پروجکشن در حین استنتاج و آموزش به صورت زیر فرمول‌بندی می‌شود:
\begin{equation}
\mathbf{W} = \mathbf{U}_{:, :r} \operatorname{diag}(\boldsymbol{\Sigma}_{1:r} \odot \boldsymbol{\gamma}) \mathbf{V}_{1:r, :}^\top + \frac{\alpha}{r} \mathbf{A} \mathbf{B}
\end{equation}
که در آن:
\begin{itemize}
    \item $\mathbf{U}_{:, :r} \in \mathbb{R}^{d_{\text{out}} \times r}$ و $\mathbf{V}_{1:r, :}^\top \in \mathbb{R}^{r \times d_{\text{in}}}$ پایه‌های متعامد مقادیر منفرد هستند که منجمد باقی می‌مانند.
    \item $\boldsymbol{\gamma} \in \mathbb{R}^r$ بردار پارامترهای آموزش‌پذیر مقیاس‌گذاری طیفی است که با ۱ مقداردهی اولیه می‌شود ($\boldsymbol{\gamma}_{\text{init}} = \mathbf{1}$).
    \item $\mathbf{A} \in \mathbb{R}^{d_{\text{out}} \times r}$ و $\mathbf{B} \in \mathbb{R}^{r \times d_{\text{in}}}$ آداپتورهای باقیمانده تصحیح‌کننده کم‌مرتبه هستند.
\end{itemize}
خروجی نگاشت به ازای تانسور ورودی $\mathbf{X}$ برابر است با:
\begin{equation}
\mathbf{Y} = \mathbf{X} \left( \mathbf{V}_{1:r, :} \operatorname{diag}(\boldsymbol{\Sigma}_{1:r} \odot \boldsymbol{\gamma}) \mathbf{U}_{:, :r}^\top \right) + \frac{\alpha}{r} \mathbf{X} \mathbf{B}^\top \mathbf{A}^\top + \mathbf{b}
\end{equation}
تانسورهای پرس‌وجو، کلید و مقدار پس از فرافکنی و اعمال کدگذاری موقعیت چرخشی (\lr{RoPE}) و هم‌ترازی سرهای \lr{GQA}، وارد لایه توجه شده و بافتار نهایی استخراج می‌گردد:
\begin{equation}
\mathbf{O} = \operatorname{softmax}\left(\frac{\mathbf{Q} \mathbf{K}^\top}{\sqrt{d}} + \mathbf{M}\right) \mathbf{V}
\end{equation}

که در آن $\sigma_i(\mathbf{W})$ بیانگر $i$-امین مقدار منفرد مرتب‌شده ماتریس به صورت نزولی است. نسبت حفظ انرژی طیفی $\eta_r = \frac{\sum_{i=1}^r \sigma_i^2}{\sum_{i=1}^{d} \sigma_i^2}$ تضمین می‌کند که با انتخاب آستانه $\tau \ge 0.90$، بیش از نود درصد از پراکندگی واریانس داده‌ها حفظ شده و افت سیگنال در انتشار گرادیان رو به عقب به حداقل ممکن کاهش می‌یابد.

\paragraph{تحلیل مجانبی و پیچیدگی اشتراک حافظه سخت‌افزاری}
از منظر پیچیدگی محاسباتی و انتقال داده‌ها در سلسومراتب حافظه (حافظه سراسری \lr{HBM} به حافظه محلی \lr{SRAM})، عملیات ضرب ماتریسی در این مکانیزم نیازمند تخصیص بافر موقت با ابعاد فاکتورگیری‌شده است. برای پردازش دنباله‌ای به طول $N$ با بعد مخفی $D$، بار محاسباتی هسته از فرم درجه دو به رابطه زیر تقلیل می‌یابد:
\begin{equation}
\label{eq:complexity_salt_attention}
\mathcal{C}_{\text{salt\_attention}} = \mathcal{O}\left( B \cdot N \cdot D \cdot r + B \cdot H \cdot N \cdot d_k \right) \ll \mathcal{O}\left( B \cdot H \cdot N^2 \cdot d_k \right)
\end{equation}
این بهبود مرتبه پیچیدگی، با کاهش چشمگیر حجم تبادل تانسورها در گذرگاه‌های ارتباطی کارت گرافیک، مانع از اشباع پهنای باند حافظه (پدیده \lr{Memory-Bound}) در فرآیند استنتاج طولانی‌مدت می‌گردد.
